\section{AION Phase 20Q --- Mission Failure Recovery Protocol}
\label{sec:aion-phase20q-mission-failure-recovery}

\subsection{Status}

Phase 20Q locks the mission failure recovery protocol.

\begin{quote}
\textbf{Status: LOCKED --- failed or blocked missions preserve useful outputs, receipts, traces, proof, and revised safe plans.}
\end{quote}

\subsection{Purpose}

AION Pilot must fail safely without losing useful work.

When a mission blocks, fails, or partially succeeds, AION must preserve trace state, proof state, receipts, and useful partial outputs.

\begin{quote}
\textbf{Failure is not data loss. A stopped mission still produces a receipt, recovery summary, and safe next plan.}
\end{quote}

\subsection{Recovery States}

The recovery protocol defines:

\begin{itemize}
  \item \texttt{failed\_recoverable}
  \item \texttt{failed\_final}
  \item \texttt{partial\_success}
  \item \texttt{blocked\_for\_safety}
\end{itemize}

\subsection{Partial Output Preservation}

Useful partial outputs MUST be preserved with:

\begin{itemize}
  \item output ID;
  \item step ID;
  \item artifact type;
  \item business container path;
  \item output hash;
  \item usability flag.
\end{itemize}

\subsection{Failure Receipt}

Every failed or blocked mission MUST emit a failure receipt containing:

\begin{itemize}
  \item mission ID;
  \item mission run ID;
  \item failure state;
  \item failure reason;
  \item stopped step ID;
  \item trace hash;
  \item proof hash;
  \item partial output hashes;
  \item receipt hashes;
  \item blocked actions;
  \item failure receipt hash.
\end{itemize}

\subsection{Recovery Summary}

A recovery summary MUST explain:

\begin{itemize}
  \item why AION stopped;
  \item which step stopped;
  \item which outputs were preserved;
  \item which actions were blocked;
  \item whether a safer revised plan is available.
\end{itemize}

\subsection{Revised Safe Plan}

After a failure, AION MAY propose a revised plan.

The revised plan MUST:

\begin{itemize}
  \item preserve safe autonomous steps;
  \item convert blocked or unsafe steps to checkpoint-required steps;
  \item disable live side effects;
  \item require human review before external, financial, legal, deployment, memory, escrow, booking, payment, dispatch, or send actions.
\end{itemize}

\subsection{Failure Invariant}

A failed mission MUST NOT lose:

\begin{itemize}
  \item trace hash;
  \item proof hash;
  \item receipt hashes;
  \item useful partial outputs;
  \item blocked action explanation;
  \item recovery hash.
\end{itemize}

\subsection{Validation}

\begin{verbatim}
python -m pytest -q \
  backend/tests/workflow_capsules/test_aion_phase20a_mission_contract_lock.py \
  backend/tests/workflow_capsules/test_aion_phase20b_mission_template_lock.py \
  backend/tests/workflow_capsules/test_aion_phase20c_mission_planner_lock.py \
  backend/tests/workflow_capsules/test_aion_phase20d_autonomy_lane_classifier_lock.py \
  backend/tests/workflow_capsules/test_aion_phase20v_kernel_guard_lock.py \
  backend/tests/workflow_capsules/test_aion_phase20e_mission_runtime_lock.py \
  backend/tests/workflow_capsules/test_aion_phase20aa_business_container_artifacts_lock.py \
  backend/tests/workflow_capsules/test_aion_phase20z_pilot_native_runtime_lock.py \
  backend/tests/workflow_capsules/test_aion_phase20w_self_repair_runtime_lock.py \
  backend/tests/workflow_capsules/test_aion_phase20x_session_vfs_lock.py \
  backend/tests/workflow_capsules/test_aion_phase20f_mission_approval_lock.py \
  backend/tests/workflow_capsules/test_aion_phase20y_canonical_payload_hashing_lock.py \
  backend/tests/workflow_capsules/test_aion_phase20r_approval_expiry_lock.py \
  backend/tests/workflow_capsules/test_aion_phase20t_mission_diff_payload_escrow_lock.py \
  backend/tests/workflow_capsules/test_aion_phase20p_mission_contract_integrity_lock.py \
  backend/tests/workflow_capsules/test_aion_phase20q_mission_failure_recovery_lock.py
\end{verbatim}

\subsection{Lock Footer}

\begin{verbatim}
Lock ID: AION-PHASE20Q-MISSION-FAILURE-RECOVERY-PROTOCOL
Status: LOCKED
Maintainer: Tessaris AI
Author: Kevin Robinson
\end{verbatim}

\subsection{Phase 20Q.1 --- Atomic Recovery Transition and Recovery Hash Closure}
\label{subsec:aion-phase20q1-atomic-recovery}

\subsubsection{Status}

Phase 20Q.1 hardens the recovery boundary with atomic transition semantics and a cryptographic recovery closure hash.

\begin{quote}
\textbf{Status: LOCKED --- recovery transitions MUST freeze execution, suspend VFS handles, snapshot state, drop process permissions, and emit deterministic closure hashes.}
\end{quote}

\subsubsection{Recovery State Transition Matrix}

To prevent unhandled exceptions or memory-exhaustion faults from dropping the runtime engine into undefined or unlogged states, the mission controller MUST implement an atomic recovery transition function $\mathcal{F}_{\mathrm{recovery}}$.

Let $\mathcal{S}_{\mathrm{current}}$ be the active execution state, and let $\mathcal{E}$ be the intercepting error or policy violation block.

The state transition is formalised as:

\begin{equation}
\mathcal{F}_{\mathrm{recovery}}
(
\mathcal{S}_{\mathrm{current}},
\mathcal{E}
)
\longrightarrow
\mathcal{S}_{\mathrm{target}}
\in
\{
\mathtt{failed\_recoverable},
\mathtt{failed\_final},
\mathtt{partial\_success},
\mathtt{blocked\_for\_safety}
\}
\end{equation}

When this transition triggers, the runtime engine MUST:

\begin{itemize}
  \item freeze all execution threads;
  \item suspend the active VFS handle;
  \item write a snapshot memory buffer to the tracking layer;
  \item drop active process permissions;
  \item preserve trace, proof, receipt and partial output state;
  \item prevent further out-of-bounds mutations.
\end{itemize}

\subsubsection{Atomic Transition Record}

The transition record MUST include:

\begin{itemize}
  \item current state;
  \item error type;
  \item target state;
  \item execution thread freeze flag;
  \item VFS suspension flag;
  \item snapshot buffer requirement flag;
  \item process permission drop flag;
  \item transition hash.
\end{itemize}

\subsubsection{Recovery Cryptographic Closure Formula}

The recovery hash $\mathcal{R}_{\mathrm{hash}}$ acts as an unalterable signature anchor sealing the system state during an abort routine.

The deterministic hash signature is defined as:

\begin{equation}
\mathcal{R}_{\mathrm{hash}}
=
\mathrm{SHA256}
(
H_{\mathrm{receipt}}
\parallel
H_{\mathrm{trace}}
\parallel
H_{\mathrm{proof}}
\parallel
\sum_i H_{\mathrm{partial\_output}_i}
\parallel
\mathtt{State}_{\mathrm{target}}
)
\end{equation}

This signature validation block MUST be pushed to the persistent logging layer before the container context relinquishes host thread priority.

\subsubsection{Validation Coverage}

Phase 20Q.1 tests prove:

\begin{itemize}
  \item atomic recovery transitions map safety violations to \texttt{blocked\_for\_safety};
  \item recoverable failures with partial outputs map to \texttt{partial\_success};
  \item recovery closure hashes are deterministic;
  \item recovery closure hashes change when target state changes;
  \item full atomic recovery includes both transition record and closure hash.
\end{itemize}

\subsubsection{Security Invariant}

\begin{quote}
\textbf{AION failure recovery is an atomic transition, not a best-effort cleanup. Once recovery starts, execution freezes, state is preserved, permissions drop, and the recovery boundary is cryptographically sealed.}
\end{quote}

\subsubsection{Lock Footer}

\begin{verbatim}
Lock ID: AION-PHASE20Q1-ATOMIC-RECOVERY-TRANSITION-CLOSURE-HASH
Status: LOCKED
Maintainer: Tessaris AI
Author: Kevin Robinson
\end{verbatim}
